%auto-ignore
%      this ensures the arxiv doesn't try to start TeXing here.
%!TEX root = super_lattice_models_draft.tex
%      prev line helps TeXShop do the right thing
\clearpage
\appendix

\section{\texorpdfstring{$\mathcal{A}_{k+1}$}{A[k+1]} models}
\label{app:TetraSym}

In this section, we describe in detail the $\mathcal{A}_{k+1}$ category, explicitly writing out its tetrahedral symbols. This category underlies the ABF height models, one of the chief examples of this paper.
The simple objects in $\mathcal{A}_{k+1}$ are labelled by $0, \tfrac12, 1, \tfrac32, \dots, \tfrac{k}{2}$.
The quantum dimensions of these objects are
\begin{align}
	d_h = \frac{ \sin\frac{\pi(2h+1)}{k+2} }{ \sin\frac{\pi}{k+2} } \,.
\end{align}
The fusion rules are
\begin{align}
	N^a_{bc} \;=\; \bigg\{ \begin{tabular}{l @{\qquad} l}
		1	&	$a+b \geq c$,\; $b+c \geq a$,\; $c+a \geq b$,\; $a+b+c \in \mathbb{Z}$,\; and $a+b+c \leq k$ \,,
	\\[0.2ex]	0	&	otherwise \,.
	\end{tabular}
\end{align}


\subsection{Tetrahedral symbols}


Let $0 \leq a,b,c,d,e,f \leq \frac{k}{2}$ be elements of the $\mathcal{A}_{k+1}$ fusion category.
One choice of the tetrahedral symbols are~\cite{kauffman1993}
\begin{align}
	\Msixj{a & b & c}{d & e & f}
   \;&=\;
	\Wsixj{a & b & c}{d & e & f} (-1)^p
	\label{eq:def_TetraSym} \,,
\end{align}
where
\begin{align}
	p &\defineas \frac{ 3(a+b+c+d+e+f)^2 - (a+d)^2 - (b+e)^2 - (c+f)^2 }{2} \,.
\end{align}
and $\{\!\begin{smallmatrix} *&*&* \\ *&*&* \end{smallmatrix}\!\}_q$ are the $q$-deformed Wigner-6j symbols (cf.\ App.~\ref{app:Wigner6j}).
Thus, the tetrahedral symbols are identical to the Wigner-6j symbols up to some signs.



\paragraph{Limiting cases.}
When one of the element is zero, the symbol reduces to
\begin{align}
	%\Msixj{a & \alpha & 0}{b & \beta & x} = \delta_{a\alpha} \delta_{b\beta} \frac{N_{ab}^x}{\sqrt{d_a d_b}} \,.
	\Msixj{\varphi & a & b}{0 & b & a} = \frac{N_{ab}^\varphi}{\sqrt{d_a d_b}} \,.
\end{align}
(Contrast with Eq.~\eqref{eq:Wigner6j_limit} for the Wigner-6j symbols.)
The expressions are also much simpler when one of the elements is $\frac12$ \cite{Lienert1992}:
\begin{subequations} \label{tetrahalf} \begin{align}
	\Msixj{\varphi & a &b}{\frac12 & b-\frac12 & a+\frac12} &= N_{ab}^\varphi N_{a+\frac12,b-\frac12}^\varphi \times (-1)^{a+b-\varphi} \sqrt{\frac{d_{\frac{\varphi+a-b}{2}} d_{\frac{\varphi+b-a-1}{2}}}{d_a d_b d_{a+\frac12} d_{b-\frac12}}} \,,
\\	\Msixj{\varphi & a &b}{\frac12 & b-\frac12 & a-\frac12} &= N_{ab}^\varphi N_{a-\frac12,b-\frac12}^\varphi \times \sqrt{\frac{d_{\frac{a+b+\varphi}{2}} d_{\frac{a+b-\varphi-1}{2}}}{d_a d_b d_{a-\frac12} d_{b-\frac12}}} \,.
\end{align}\end{subequations}


\subsection{\texorpdfstring{Wigner-6j symbols}{Wigner-6j symbols}}
\label{app:Wigner6j}
The $q$-deformed Wigner-6j symbols may be computed via the Racah formula \cite{Reshetikhin1988}.
\begin{align}\begin{split}
	\Wsixj{j_1 & j_2 & j_3}{j_4 & j_5 & j_6}
	&= \Delta(j_1,j_2,j_3) \, \Delta(j_1,j_5,j_6) \, \Delta(j_4,j_2,j_6) \, \Delta(j_4,j_5,j_3)
	\\	&\quad
		\times \sum_z \begin{pmatrix}
				\frac{ (-1)^z \; [z+1]! }{ [z-j_1-j_2-j_3]! \ [z-j_1-j_5-j_6]! \ [z-j_4-j_2-j_6]! \ [z-j_4-j_5-j_3]! }
			\\[1ex]	\times
				\frac{1}{ [j_1+j_2+j_4+j_5-z]! \ [j_2+j_3+j_5+j_6-z]! \ [j_3+j_1+j_6+j_4-z]! }
			\end{pmatrix} ,
	\label{eq:def_W6j}
\end{split}\end{align}
where
\begin{align}\begin{split}
	q &\defineas \exp\tfrac{2 \pi i}{k+2} \,,
\\	[m] &\defineas \frac{q^{m/2} - q^{-m/2}}{q^{1/2} - q^{-1/2}}
		= \frac{\sin \frac{m\pi}{k+2}}{\sin \frac{\pi}{k+2}} \,,
\\	[n]! &\defineas \left\{ \! \begin{array}{cl}
			\prod_{m=1}^n [m]  &  n > 0 \,,
		\\	1  &  n = 0 \,,
		\\	\infty  &  n < 0 \,,
		\end{array} \right.
\\	\Delta(j_1, j_2, j_3) &\defineas \left\{ \! \begin{array}{cl}
		\sqrt{\dfrac{ [j_1+j_2-j_3]! \ [j_3+j_1-j_2]! \ [j_2+j_3-j_1]! }{ [j_1+j_2+j_3+1]! }}	&	N^{j_1}_{j_2 j_3} = 1 \,, \\[3ex]
		0	&	\text{otherwise} \,, \end{array} \right.
\end{split}\end{align}
The summation $\sum_z$ is over integers $z$ such that the summand is well-behaved (i.e., no terms being zero or infinity).
Since $[k+2] = 0$, thus $[n]!$ vanishes for $n \geq k+2$.
Also note that $[n]$ yields all the quantum dimensions as $d_x = [2x+1]$.

The Wigner-6j symbols is invariant under the same permutations as those of the tetrahedral symbols, i.e., Eqs.~\eqref{eq:tetrasym_permute_cols}, \eqref{eq:tetrasym_swap_rows}.
However, they have different limiting cases.
\begin{align}
	\Wsixj{a & a & 0}{b & b & c} = \frac{N^a_{bc} (-1)^{a+b+c}}{\sqrt{d_a d_b}} .
	\label{eq:Wigner6j_limit}
\end{align}
We also have
\begin{align}
	\Wsixj{a & b & c}{\frac{k}{2}-d & \frac{k}{2}-e & \frac{k}{2}-f}
		&= (-1)^{k+a+b+c}
		\Wsixj{a & b & c}{d & e & f} .
\end{align}
It follows that
\begin{align}
	\renewcommand{\arraystretch}{1.1}
	\Wsixj{\frac{k}{2}-a & \frac{k}{2}-b & c}{\frac{k}{2}-d & \frac{k}{2}-e & \quad f\quad}
		&= (-1)^{k+a+b+d+e}
		\Wsixj{a & b & c}{d & e & f} .
\end{align}
The associated $F$-symbols are defined as
\begin{align}
	\big[F_e^{abd}\big]_{c,f} = (-1)^{a+b+d+e} \Wsixj{a&b&c}{d&e&f} \sqrt{d_c d_f} .
\end{align}
Hence, the Racah-Elliot-5 relation also appears different.
\begin{align}
	\sum_x d_x (-1)^{S}
		\Wsixj{a & b & x}{c & d & g} \Wsixj{c & d & x}{e & f & h} \Wsixj{e & f & x}{b & a & j}
	&= \Wsixj{g & h & j}{e & a & d} \Wsixj{g & h & j}{f & b & c} ,
\end{align}
with the sign determined by $S = a+b+c+d+e+f+g+h+j+x$.





\comment{
\section{Properties of tetrahedral symbols}
\label{app:tetra}

\paragraph{Orthogonality relations.} The tetrahedral symbols satisfies the \emph{Racah-Elliot} orthgonality relations.
\\
Racah-Elliot-3:
\begin{align}
	\sum_x d_x \Msixj{a & b & m}{c & d & x} \Msixj{a & b & n}{c & d & x}
	&= \frac{\delta_{mn}}{d_m} N^m_{ab} N^m_{cd} .
	\label{eq:RE3}
\end{align}
Racah-Elliot-5:
\begin{align}
	\sum_x d_x \Msixj{a & b & x}{c & d & g} \Msixj{d & c & x}{f & e & h} \Msixj{e & f & x}{b & a & j}
	&= \Msixj{g & h & j}{e & a & d} \Msixj{g & h & j}{f & b & c} .
	\label{eq:RE5}
\end{align}
Or equivalently
\begin{align}
	\sum_x d_x \Msixj{\beta & \gamma & g}{\rho & \nu & x} \Msixj{\gamma & \alpha & h}{\mu & \rho & x} \Msixj{\alpha & \beta & j}{\nu & \mu & x}
	&= \Msixj{g & h & j}{\alpha & \beta & \gamma} \Msixj{g & h & j}{\mu & \nu & \rho} .
\end{align}
\paragraph{$F$-symbols.} The $F$-symbols for the equivalent tensor category are
\begin{align}
	\big[ F_d^{abc} \big]_{xy} &= \sqrt{d_{xy}} \Msixj{a & b & x}{c & d & y} .
\end{align}
Thus Racah-Elliot-3 implies unitarity of the tensor category.

\begin{align}
	\sum_x
		\big[ F_a^{bcd} \big]_{g,x} \big[ F_j^{bxe} \big]_{a,f} \big[ F_f^{cde} \big]_{x,h}
	=	\big[ F_j^{gde} \big]_{a,h} \big[ F_j^{bch} \big]_{g,f} \;.
\end{align}
Racah-Elliot-4 is equivalent to the \emph{hexagon equation} while Racah-Elliot-5 is equivalent to the \emph{pentagon equation}.

$R$-symbols.
\begin{eqnarray}
R_c^{ab} &=& e^{i\pi (h_c-h_a-h_b)}\\
h_a &=& a^2-\frac{a(a+1)}{k+2}
\end{eqnarray}
satisfy Bonderson 2.78


}


%%%%%%%%%%%%%%%%%%%%%%%%%%%%%%%%%%%%%%%%%%%%%%%%%%%%%%%%%%%%%%%%%%%%%%%%%%%%%%%%
\section{6j symbols for fusion categories with non-self-conjugate objects}
\label{app:notselfdual}

A $\mathbb{C}$-linear, semi-simple, rigid monoidal categories is really just a a fancy name for a fusion category.
And since everyone already knows what a spherical, unitary fusion category is, we just have to explain the remaining qualifiers.

The point of this section is to speculate on cases where $x \neq \bar{x}$ for some $x$.
It is also important to distinguish between upper/lower indices in the fusion symbol
\begin{align}
	N^{c}_{ab} = N^{b}_{c\bar{a}} = N^{\bar{a}}_{b\bar{c}} = N^{\bar{c}}_{\bar{a}\bar{b}} = N^{\bar{b}}_{\bar{c}a} = N^{a}_{\bar{b}c} .
\end{align}
(Note $N^\ast_{\ast\ast}$ is invariant under exchanging the two lower indicies: $N^c_{ab} = N^c_{ba}$.)
For each $a$, its quantum dimension $d_a$ is the dominant eigenvalue of the matrix $N^\ast_{a\ast}$.
They quantum dimensions satisfy
\begin{align}
	d_0 &= 1 \,,
&	d_a &\geq 1 \,,
&	d_{\bar{a}} &= d_a \,,
&	\sum_x N^x_{ab} d_x &= d_a d_b \,.
\end{align}
We require the theory to be \emph{multiplicity-free}, that is, $N^c_{ab} \in \{0, 1\}$ be no larger than 1; thus no additional degrees of freedom lie at the graph vertices.
The $A$ and $B$-symbols (in absence of fusion multiplicities) are defined as
\begin{subequations}\begin{align}
	A^{ab}_c & \;\defineas\; \sqrt{\frac{d_ad_b}{d_c}} \big[ F_b^{\bar{a}ab} \big]^\ast_{0,c} \;,
\\	B^{ab}_c & \;\defineas\; \sqrt{\frac{d_ad_b}{d_c}} \big[ F_c^{c\bar{b}b} \big]_{a,0} \;.
\end{align}\end{subequations}
Note that the Frobenius-Schur indicators $\varkappa_a = A_0^{a\bar{a}} = B_a^{0a}$ are special cases of these symbols.

We establish the criteria for which the tetrahedral tetrahedral are defined.
First we must fix the gauge, i.e., choose a set of basis vectors $\psi_c^{ab} \in V_c^{ab}$, which in turns fixes the $F$-symbols.
\begin{align}
	\fbox{\text{The tetrahedral symbols are defined if $N_{ab}^c = 1 \Longrightarrow A_c^{ab} = B_c^{ab} = +1$ for all $a,b,c$.}}
\end{align}
Note that the symbols are gauge-dependent.  Hence some categories may admit one or multiple tetrahedral symbols, while some may admit none.
(For example, we require that all Frobenius-Schur indicator must be $+1$, which rules out $\SU2_k$, since $\varkappa_x = (-1)^{2x}$ in the theory.)



\subsection*{Properties of tetrahedral symbols.}

First we define the tetrahedral symbols
\begin{align}
	\Msixj{a & b & c}{d & e & f}
	& \;\defineas\;
	\frac{1}{\sqrt{d_{abcdef}}} \tetradiagram .
\end{align}
Recall that
\begin{align}
	\big[ F_d^{abc} \big]_{x,y} = \Braket{ \FmoveRsmall | \FmoveLsmall } = \frac{1}{\sqrt{d_{abcd}}} \FmoveC \;.
	\label{eq:def_Tetra22}
\end{align}
Therefore, the tetrahedral symbols are related to the $F$-symbols via
\begin{align}
	\big[ F_d^{abc} \big]_{x,y} &= \sqrt{d_{xy}} \Msixj{a & b & \bar{x}}{c & d & y} .
\end{align}

\paragraph{Properties.}
Each vertex imposes a constraint of the fusion rules:
\begin{align}
	N^{\bar{a}}_{bc} N^d_{ec} N^e_{fa} N^f_{db} = 0
	\;\Rightarrow\; \Msixj{a & b & c}{d & e & f} = 0 \,.
\end{align}
The symbols are invariant under cyclic permutation of the columns.
\begin{align}
	\Msixj{a & b & c}{d & e & f} = \Msixj{b & c & a}{e & f & d} = \Msixj{c & a & b}{f & d & e} .
\end{align}
When any two columns are interchanged, we take the dual of the top row and the conjugate the entire symbol.
\begin{align}
\label{colexchange}
	\begin{array}{ccccccc @{} l}
		\Msixj{a & b & c}{d & e & f}
	&	=	&	\Msixj{\bar{a} & \bar{c} & \bar{b}}{d & f & e}^\ast
	&	=	&	\Msixj{\bar{c} & \bar{b} & \bar{a}}{f & e & d}^\ast
	&	=	&	\Msixj{\bar{b} & \bar{a} & \bar{c}}{e & d & f}^\ast
	&	.
	\\[-1.6ex]
	&&	\hspace{3ex} \underbracket{\hspace{4ex}}_{\text{exchange}} \hspace{0.7ex}
	&&	\underbracket{\hspace{7ex}}_{\text{exchange}} \hspace{0.7ex}
	&&	\underbracket{\hspace{4ex}}_{\text{exchange}} \hspace{4ex}
	\end{array}
\end{align}
We may also exchange the rows among any two columns, at the cost of dualizing some entries.
\begin{align}
	\begin{array}{ccccccc @{} l}
		\Msixj{a & b & c}{d & e & f}
	&	=	&	\Msixj{\bar{a} & e & \bar{f}}{\bar{d} & b & \bar{c}}
	&	=	&	\Msixj{\bar{d} & \bar{b} & f}{\bar{a} & \bar{e} & c}
	&	=	&	\Msixj{d & \bar{e} & \bar{c}}{a & \bar{b} & \bar{f}}
	&	.
	\\[-1.6ex]
	&&	\hspace{3ex}\underbracket{\hspace{4ex}}_{\text{swap rows}}
	&&	\underbracket{\hspace{7ex}}_{\text{swap rows}}
	&&	\underbracket{\hspace{4ex}}_{\text{swap rows}}\hspace{3ex}
	\end{array}
\end{align}
Finally, we also have the limiting case
\begin{align}
	\Msixj{\bar{a} & a & 0}{b & b & c}
	= \Msixj{a & b & \bar{c}}{\bar{b} & a & 0}
	= \frac{N^c_{ab}}{\sqrt{d_ad_b}} \,.
\end{align}


\paragraph{Orthogonality relations.}
The unitarity of the $F$-symbols imply
\begin{align}
	\sum_y d_y \Msixj{a & b & x}{c & d & y} \Msixj{a & b & x'}{c & d & y}^\ast
	= \frac{ \delta_{xx'} N^{\bar{x}}_{ab} N^{x}_{c\bar{d}} }{d_x} \,,
\end{align}
analogous to Racah-Elliot-3.

For the analogous equation to Racah-Elliot-5, consider the following diagram.
\begin{align}
	\circlespentagon
\end{align}
After evaluating the diagram two different ways,%
	\footnote{%
		The left hand side comes from applying three $F$-moves to eliminate $g$, $h$, and $j$,
			resulting in coefficients $\big[F_\gamma^{\rho\bar\nu\beta}\big]_{\bar{g},x}$, etc.
		A useful identity for the right hand side is $\vertexabcC = \sqrt{d_{def}} \Msixj{a & b & c}{d & e & f}^\ast \vertexabc$.}
	eliminating the common factor $\sqrt{d_{\alpha\beta\gamma ghj \mu\nu\rho}}$ from both sides, we get the pentagon equation:
\begin{align}
	\sum_x d_x
		\Msixj{\rho & \bar\nu & g}{\beta & \gamma & x}
		\Msixj{\mu & \bar\rho & h}{\gamma & \alpha & x}
		\Msixj{\nu & \bar\mu & j}{\alpha & \beta & x}
	=	\Msixj{g & h & j}{\alpha & \beta & \gamma}
		\Msixj{g & h & j}{\mu & \nu & \rho}^\ast .
	\label{eq:super_pentagon}
\end{align}


\paragraph{Other relations.}
\begin{align}
	\sum_f d_f \Msixj{a & \bar{a} & c}{b & \bar{b} & f} &= \delta_{c,0} \sqrt{d_{ab}} \,,
\\	\sum_{x,\mu,\nu,\rho} d_x d_\mu d_\nu d_\rho \Msixj{\rho & \bar\nu & g}{\beta & \gamma & x} 
		\Msixj{\mu & \bar\rho & h}{\gamma & \alpha & x}
		\Msixj{\nu & \bar\mu & j}{\alpha & \beta & x}
		\Msixj{g & h & j}{\mu & \nu & \rho}
	&=	\mathscr{D}^2 \Msixj{g & h & j}{\alpha & \beta & \gamma} .
\end{align}





\section{Tambara-Yamagami fusion categories and tetrahedral symbols}
\label{app:Potts}

In this appendix we review the Tambara-Yamagami fusion categories and write down their 6j-symbols~\cite{Tambara1998}.

Let $A$ be an Abelian group. 
The fusion category has one object for every group with quantum dimension equal to 1, and a single non-Abelian object $X$ with quantum dimensions $d_X = \sqrt{|A|}$, where $|A|$ is the number of group elements in $A$.
The identity element $\I$ corresponds to $0 \in A$.
The fusion rules are given by,
\begin{align}
	a \otimes b = (a{+}b) \,, \qquad a \otimes X = X \otimes a = X \,, \qquad X \otimes X = \; \bigoplus_{a \in A} a \,,
\end{align}
where $a,b$ are elements of $A$ and $a+b$ denotes group composition in $A$.
We also have that $\bar{X} = X$ and $\bar{a} = -a$. 
The fusion category is $\mathbb{Z}_2$ graded with $\mcc_0 = \{a \in A \}$ and $\mcc_1 = \{ X \}$.

The 6j-symbols depend on a symmetric non-degenerate bi-character $\chi$,
That is, $\chi: A \times A \rightarrow \mathbb{C}^\times$ with $\chi(a+b,c) = \chi(a,c)\chi(b,c)$, $\chi(a,b+c) = \chi(a,b)\chi(a,c)$, $\chi(a,b) = \chi(b,a)$, and $\det \chi \neq 0$. 
We also work in a gauge where $\chi(0,a) = \chi(a,0) = 1$, where 0 is the identity element of $A$.
The 6j-symbols are given by 
\begin{align}
\label{ATY6j}
	\Msixj{b{-}c & c{-}a & a{-}b}{a & b & c} &= 1 \,,
&	\Msixj{X & X & a}{X & X & b} &= \frac{ \chi(a,b)}{d_X} \,,
&	\Msixj{a & b & c}{X & X & X} &= \frac{N_{bc}^{\bar{a}}}{\sqrt{d_X}} \,.
\end{align}
The remaining 6j-symbols can be inferred from symmetry.



\subsection*{Example: \texorpdfstring{$\mathbb{Z}_3$}{Z3} Tambara-Yamagami}

The simple objects are $0$, $1$, $2$, and $X$.
The non-trivial fusion symbols are
\begin{align}
	N^0_{XX} = N^1_{XX} = N^2_{XX} = 1 \,,
	\quad	N^X_{X1} = N^X_{2X} = 1 \,,
	\quad	N^0_{12} = 1 \,,
	\quad	N^1_{22} = N^2_{11} = 1 \,.
\end{align}
Therefore, $\bar{1} = 2$ and $\bar{X} = X$.
The quantum dimensions are
\begin{align}
	d_0 = d_1 = d_2 = 1 \,,
	\quad	d_X = \sqrt{3} \,.
\end{align}
The only non-Abelian anyon is $X$, and its $F$-symbol is given by
\begin{align}
	\big[ F_X^{XXX} \big]_{a,b} = \frac{1}{\sqrt{3}}
			\begin{pmatrix} 1 & 1 & 1 \\ 1 & \omega & \omega^2 \\ 1 & \omega^2 & \omega \end{pmatrix}_{a,b}
		&=	\frac{1}{\sqrt{3}} \omega^{ab} \,,
	&	\text{where } \omega &= e^{\frac{2\pi i}{3}} \,.
\end{align}
Therefore, the tetrahedral symbols are given by
\begin{align}
	\Msixj{X & X & a}{X & X & b} &= \frac{\omega^{-ab}}{\sqrt{3}} ,
&	\Msixj{1 & 1 & 1}{X & X & X} &= \Msixj{2 & 2 & 2}{X & X & X} = \frac{1}{\sqrt[4]{3}} .
\end{align}
The remaining symbols may be inferred from symmetry.




